%%
%% Author: shoupu_wan
%% 8/6/18
%%

\chapter{Contents at a Glance}\label{sec:structureOfTheBook}

The book is divided into three parts.
Most of the material covered in this book are so basic that they may almost appear simplistic to
you.
It is exactly these basic building blocks that more useful things are based on.

In \fullref{prt:flowControl}, we lay the foundations of concise and readable
modular designs and implementations.
In \fullref{ch:reusableFunctions} we touch base on well-structured, modular, and reusable
functions.
Then we spend \fullref{ch:organizeYourIfElse} to discuss the meat of readable and succinct flow of
control pros and cons.
We conclude this part of the book by \fullref{ch:loops} and \fullref{ch:recursion} where some
insights of loops and recursions are presented. We will talk about the ``Loop and a half''
problem and possible solutions.

I dedicate \fullref{prt:dataStructure} to explain how I settled
with a handful of usage patterns that can also help you write reliable, maintainable code.
With a plenty of examples, we will take look at the usage
patterns of common data structures such as maps, linked lists, binary trees, stacks, and queues
(\fullref{ch:ds}).
For each of these data structures, I assume you already have a solid understanding of what they
are.
\fullref{ch:knowTheLanguageFeatures} talks about specific language features that may help
you write concise and readable code.
But we will maintain the practical purpose of readability instead of showing off tricky, enigmatic
esoteric `language features'.


With the foundations laid, We will go at length on a few advanced topics in \fullref{prt:advanced}.
We will start with \fullref{ch:loop-thinning} where we explain how to simplify nested loops by
introducing a technology dubbed as ``loop thinning''\index{Loop Thinning\emph{Loop Thinning}}.
In the next \fullref{ch:sentinels}, we will talk about sentinels and their application to reduce
redundant and / or duplicate code.
Another topic I'd like to elaborate on is dynamic programming and this will be done in \fullref{ch:dp}.
I will show you how putting on the right mindset can help you make friend with dynamic programming.
I will also explain the two common approaches to crack a complex dynamic programming problem.
Furthermore, I go into details to explain a multi-variate dynamic programming techniques for
solving complicated dynamic programming problems.
Searching through the literature, I haven't found obvious sources lecturing on iterative approach to
object-oriented design (OOD) during implementation stage of software development cycles.
When a conventional book talks about OOD, it is often done in an isolated way.
Seldom any book has touched the interaction between designs and implementations.
This is exactly my focus in \fullref{ch:ood}.
In a section, I will attempt to explain how an iterative approach can bring a near optimal
design and implementation by taking advantage of the interplay between the two.
Finally, I will dedicate a section to explain a methodology that uses variational principle to
approach an optimal design and implementation of a complex piece of software.
